\subsection{Supersingular correspondences: a larger proved class}

The general assertion in B5 concerns a finite \'etale cover whose
geometric Galois closure has order prime to $p$. This condition is
stronger than requiring the degree of a particular non-Galois cover
to be prime to $p$. It is retained throughout the constructions below.

\begin{claim}{B.LIT1}{Proved: placing finitely many values over prime-to-$p$ torsion}{B5}
Let $E/\overline{\F}_p$ be an elliptic curve, let
$\pi:E\to E/\{\pm1\}\simeq\mathbf P^1$ be its quotient map, and
let $S\subset\mathbf P^1(\overline{\F}_p)$ be finite. There exists
$\phi\in\mathrm{PGL}_2(\overline{\F}_p)$ such that every point of
$\phi(S)$ is the image under $\pi$ of a prime-to-$p$ torsion point
of $E$. The points of $\phi(S)$ may also be required to avoid the
branch locus of $\pi$.
\end{claim}
\begin{proof}
Choose coordinates putting both $S=\{t_1,\ldots,t_N\}$ and the
branch locus $B$ of $\pi$ in $\mathbf A^1$. This is permissible even
in characteristic two; the quotient map is then separable of degree
two, although ramification may be wild. Choose $c\ne0$ outside the
finite set determined by
\[
 c(t_i-t_j)=b-b'\qquad(i\ne j,\ b,b'\in B).
\]
Thus the finite sets $B-ct_i$ are pairwise disjoint. Descend all this
data to a finite field $k=\F_q$.

If $E$ is supersingular, put $a=0$ and $v=\operatorname{id}_E$.
If $E$ is ordinary, its $p$-adic Tate module has rank one. After
enlarging $k$, write the Frobenius action as multiplication by
$\alpha\in1+p\Z_p$, and put $a=v_p(\alpha-1)<\infty$.
For every $n$ prime to $p$,
\[
 v_p(\alpha^n-1)=a,
\]
because $(\alpha^n-1)/(\alpha-1)\equiv n\pmod p$.
Thus the $p$-primary subgroup of $E(\F_{q^n})$ is killed by $p^a$.
Let $v:E'=E^{(p^a)}\to E$ be the iterated Verschiebung. It is an
\'etale isogeny of degree $p^a$ in the ordinary case, and
\[
 v(E'(\F_{q^n}))=[p^a]E(\F_{q^n})
                 =E(\F_{q^n})\{p'\},\qquad p\nmid n.
\]
Here relative Frobenius is a bijection on points over a perfect
finite field and $v\circ F^a=[p^a]$. In the supersingular case all
geometric torsion already has order prime to $p$.

The composite $h=\pi\circ v:E'\to\mathbf P^1$ is Galois with
group $G=\ker(v)\rtimes\{\pm1\}$ and branch locus $B$.
Its branch inertia groups normally generate $G$: the quotient by
their normal closure would be a connected finite \'etale cover of
$\mathbf P^1$, which must have degree one by Riemann--Hurwitz.
Consider the $N$ covers of the $s$-line given by
\[
 h(R_i)=s+ct_i.
\]
Each is connected and has branch locus $B-ct_i$. Their compositum
has geometric Galois group $H\subseteq G^N$ projecting onto each
factor. At a branch point belonging only to the $i$th factor, its
inertia subgroup is supported in that factor and projects onto the
corresponding inertia of $G$. The other covers are unramified there.
Conjugating by $H$ and using the surjection to the $i$th factor shows
that $H$ contains the normal closure of every such inertia in that
factor. These normal closures generate $G$. Consequently $H$ contains
every factor and equals $G^N$. This also proves geometric
irreducibility of the normalized fibre product; it does not assume
linear disjointness without checking it.

Let $Y/k$ be its smooth proper model. Fix a prime $l\ne p$ and take
$n=l^m$. The Weil bound makes $|Y(\F_{q^n})|$ tend to infinity.
Only finitely many geometric points lie above the branch points or
infinity, so for sufficiently large $m$ there is a point outside that
set. It supplies $s,R_1,\ldots,R_N$ over $\F_{q^n}$ with
$s+ct_i=\pi(v(R_i))$. Each $v(R_i)$ has prime-to-$p$ order by the
previous paragraph. The transformation $t\mapsto s+ct$, composed
with the initial coordinate change of $S$, is the desired $\phi$.
\end{proof}

\paragraph{A finite point-count bound.}
The curve in the proof has degree $d^N$ over the $s$-line, where
$d=2p^a$. Since each individual cover has genus one and their branch
loci are disjoint, the different formula gives
$g(Y)=1+(N-1)d^N$, also when $p=2$. Thus the final existence step
can be made effective by choosing an $l$-power extension for which
the lower Weil bound exceeds the finite number of excluded geometric
points. The $p$-power isogeny in this auxiliary proof is not one of
the \'etale covers of the original curve in the next construction.

\begin{claim}{B.LIT2}{Proved: every odd-characteristic hyperelliptic source}{B5}
Let $p$ be odd and let $C/\overline{\F}_p$ be hyperelliptic of genus
at least two. For every hyperelliptic curve $H/\overline{\F}_p$ there
is a connected finite \'etale cover $D\to C$ with prime-to-$p$
geometric Galois closure and a nonconstant map $D\to H$.
In particular B5 holds for every such $C$, including characteristic
three, with a supersingular target of genus at least two.
\end{claim}
\begin{proof}
Let $C\to\mathbf P^1$ be the degree-two map, with branch set $T$
of cardinality $2g(C)+2$. Choose four points of $T$ and let
$E\to\mathbf P^1$ be the elliptic double cover branched there.
The normalization $C_0$ of $C\times_{\mathbf P^1}E$ is connected:
the two quadratic function fields are different because their
branch sets are different. By tame ramification cancellation,
$C_0\to C$ is \'etale of degree two. The map $C_0\to E$ is a
double cover ramified above the points lying over the remaining
members of $T$. Translate the group law of $E$ so that one such
branch point is its origin. Then the entire fibre of $C_0\to E$
over zero has ramification index two.

Choose the quotient map $\pi:E\to E/\{\pm1\}$ and write $S$
for the branch set of a hyperelliptic map $H\to\mathbf P^1$.
By B.LIT1, after changing the coordinate on that target, its branch
points are $\pi(P_i)$ for finitely many prime-to-$p$ torsion points
$P_i\in E$. Choose an integer $m$ prime to $p$ killing all $P_i$.
Pull $C_0\to E$ back along $[m]:E\to E$ and normalize, obtaining
$C_1$. The map $C_1\to C_0$ is finite \'etale with group $E[m]$.
The double cover $C_1\to E$ is still connected, since its nonempty
branch divisor is the inverse image of the original branch divisor.
Above every point of $E[m]$ its ramification index is two.

Consequently every point of $C_1$ above a branch value $\pi(P_i)$
has even ramification index in the composite $C_1\to E\to\mathbf P^1$:
the two inverse images under $\pi$ are $P_i,-P_i\in E[m]$.
The normalization of $C_1\times_{\mathbf P^1}H$ is therefore
\'etale over $C_1$, by Abhyankar's lemma. Any connected component
$D$ dominates both $C_1$ and $H$, and has degree one or two over
$C_1$. The tower $D\to C_1\to C_0\to C$ consists of finite
\'etale Galois steps of prime-to-$p$ orders. Its Galois closure has
prime-to-$p$ order: inductively its group is a subgroup of an
extension of the preceding Galois group by a product of conjugate
step groups, equivalently of the corresponding finite wreath
product. This verifies the stronger Galois-closure condition.

For an explicit supersingular target choose an odd prime power
$Q=p^b\ge5$ and let
\[
 H_Q:\quad y^2=x^Q-x.
\]
Its genus is $(Q-1)/2$. For $x\in\F_{Q^2}\setminus\F_Q$, the
nonzero value $x^Q-x$ has trace zero. All such values have the same
quadratic character
$\epsilon=(-1)^{(Q+1)/2}$ in $\F_{Q^2}$: writing a trace-zero
element as $a$, one has $a^{Q-1}=-1$ and hence
$a^{(Q^2-1)/2}=(-1)^{(Q+1)/2}$; elements of $\F_Q^\times$ are
squares in $\F_{Q^2}$. Therefore
\[
 |H_Q(\F_{Q^2})|=Q^2+1+\epsilon(Q^2-Q).
\]
This attains one of the Weil bounds. All Frobenius eigenvalues
over $\F_{Q^2}$ are consequently $-\epsilon Q$, so all Newton
slopes are $1/2$. Thus $H_Q$ is supersingular. For $p=3$, take
$Q=9$, yielding genus four; for $p\ge5$, take $Q=p$.
\end{proof}

The first double-cover construction is the one used in
Bogomolov--Tschinkel's hyperelliptic argument. The extra assertion
proved above is that the finite torsion positioning can be arranged
with prime-to-$p$ torsion, which keeps the final covering tower in
the required formation. We make no priority claim for this
strengthening. The ramification cancellation is the precise
one-dimensional form of SGA1 XIII Proposition~5.2 and Poonen
Lemma~2.1.\citep{LA_BogomolovTschinkel2002,LA_Poonen2005,LA_SGA1}

\begin{claim}{B.LIT3}{Proved: a characteristic-two ramification criterion}{B5}
Let $p=2$. Suppose a curve $C/\overline{\F}_2$ admits a finite
separable map to $\mathbf P^1$ with three distinct fibres in which
every ramification index is divisible by three. Then a connected
finite \'etale cover of $C$ of degree one or three maps nonconstantly
to a supersingular elliptic curve. In particular this applies to
every cyclic degree-three cover of $\mathbf P^1$ of genus at least
two.
\end{claim}
\begin{proof}
Move the three values to $0,1,\infty$ and use the cyclic cover
$E:y^3=x(x-1)$. Its smooth projective model is elliptic, branched
with index three at those values. Normalized base change along
$C\to\mathbf P^1$ is \'etale by Abhyankar's lemma, and any
connected component has cyclic degree one or three over $C$.
After interchanging the coordinates, the elliptic curve is
$z^2+z=t^3$. It has nine points over $\F_4$: for each of the four
values of $t$, the right side has trace zero to $\F_2$, so there
are two values of $z$, in addition to the point at infinity.
Thus its two $\F_4$-Frobenius eigenvalues are $-2$ and it is
supersingular. A cyclic cubic cover has uniform ramification
index three over each branch point, and has at least three such
points when its genus is at least two.
\end{proof}

\begin{claim}{B.LIT4}{Proved: the first three low-genus Galois cases}{B5}
Suppose $p$ is odd, $C/\overline{\F}_p$ has genus at least two,
and $G\subset\Aut(C)$ has order prime to $p$. Put $Y=C/G$ and
suppose $g(Y)\le2$. If at least one of the following holds, then
$C$ has a connected finite \'etale cover with prime-to-$p$ Galois
closure mapping to a supersingular curve of genus at least two:
\begin{enumerate}
\item $g(Y)=1$ or $2$;
\item $G$ is solvable;
\item two distinct branch values of $C\to Y$ have ramification
indices with a common prime factor.
\end{enumerate}
In particular this proves B5 for every tamely superelliptic curve
in every odd characteristic, including characteristic three.
\end{claim}
\begin{proof}
We verify the extra prime-to-$p$ condition in Poonen's
Theorem~1.7, cases~(1)--(3), rather than infer it from the word
``unramified.'' His theorem assumes precisely a prime-to-$p$
Galois action, genus at least two for the source, and quotient
genus at most two, with one of the listed extra hypotheses.
The proof below follows his \S\S3.1--3.3 and records the orders
of the covering groups.\citep[Theorem~1.7 and \S3]{LA_Poonen2005}

First suppose $Y$ has genus two. It is already hyperelliptic,
so no \'etale covering is needed to dominate a hyperelliptic curve.
Suppose instead that $Y=E$ is elliptic. The map $C\to E$ is
ramified, since $g(C)>1$. Translate a branch value to the origin,
and choose a prime $e\ne p$ dividing its ramification index.
The Galois hypothesis ensures that this divisibility holds at
every point above the origin. Pull back along $[l]:E\to E$ for
an odd prime $l\ge5$ distinct from $p$, and choose a connected
component $C'$. Then $C'\to C$ is finite \'etale Galois with
group a subgroup of $E[l]$. Every point of $C'$ above $E[l]$
has ramification index divisible by $e$ in $C'\to E$.

Under $E\to E/\{\pm1\}=\mathbf P^1$, the nonzero $l$-torsion
points give $(l^2-1)/2$ distinct values $a_i$. If $e\ge3$, take
the cyclic cover
\[
 H:\quad z^e=\frac{(x-a_1)(x-a_2)}{(x-a_3)(x-a_4)}.
\]
It has genus $e-1\ge2$, and its $z$-map has degree two, so it
is hyperelliptic. If $e=2$, use three factors in the numerator
and three in the denominator, giving a genus-two hyperelliptic
curve instead. At every branch value of $H\to\mathbf P^1$,
all ramification indices of $C'\to\mathbf P^1$ are divisible
by $e$. A connected component of the normalized pullback of
$H$ is therefore \'etale over $C'$, with cyclic group of order
one or $e$, and dominates $H$.

Now let $Y=\mathbf P^1$ and $G$ be solvable. Induct on $|G|$.
For a nontrivial proper normal subgroup $N$, if $g(C/N)\ge2$
apply induction to $C/N\to\mathbf P^1$ and pull its \'etale
cover back to $C$. If $g(C/N)=1$, apply the elliptic case to
$C\to C/N$. If $g(C/N)=0$, apply induction to $C\to C/N$.
It remains to treat $G$ cyclic of prime order $e\ne p$.
For $e=2$, $C$ itself is hyperelliptic. For $e\ge3$, pull
back along a cyclic degree-$e$ map $\mathbf P^1\to\mathbf P^1$
branched at two branch values of $C$. The two degree-$e$
extensions are different (the original curve has positive genus),
so the pullback is connected and gives an \'etale cyclic
degree-$e$ cover $C'\to C$. Its new cyclic map to $\mathbf P^1$
has at least four branch values, since its genus has increased.
Using four of these values in the preceding displayed equation
again gives an \'etale cyclic pullback dominating a hyperelliptic
curve.

Finally suppose two branch indices have a common prime divisor
$e$. Pull back along a cyclic degree-$e$ map of projective lines
branched at those values. A third branch value has at least two
inverse images; a prime factor $e'$ of its ramification index
therefore occurs at two new branch values. Repeat this construction.
Every step on the source is cyclic \'etale of order dividing the
chosen prime, hence prime to $p$. Every new map to the projective
line is Galois of degree at most $|G|$, with ramification indices
dividing $|G|$. The degrees of the composite maps of projective
lines tend to infinity, forcing the degrees of the covers of $C$
and hence their genera to tend to infinity. Riemann--Hurwitz then
forces the number of branch values to tend to infinity. There are
only finitely many possible indices, so at some stage at least six
branch values have a common index $a>1$. Choose a prime $e\mid a$
and a cyclic degree-$e$ cover of $\mathbf P^1$ branched at these
six points, using three exponents $1$ and three exponents $-1$.
It has genus $2(e-1)\ge2$. Its normalized pullback is \'etale
over the stage curve, and its solvable cyclic group puts the target
in the preceding case.

Thus in every case a hyperelliptic curve is dominated after a
finite tower of \'etale Galois prime-to-$p$ covers and their base
changes. A connected base change of such a Galois cover has a
subgroup of the original group; finite composition has prime-to-$p$
Galois closure by the wreath-product argument in B.LIT2. Composing
with B.LIT2 (and pulling its cover back) gives the claimed
supersingular target.
\end{proof}

\paragraph{The modular fourth case, with its group condition.}
Poonen's fourth condition asks for three branch indices divisible
by $2,3,l$, respectively, where $l\le89$ is prime or
$l\in\{101,103,107,131,167,191\}$. For $l=2,3$ the preceding
common-factor case applies. For $l\ge5$ his proof uses
$X(l)\to\mathbf P^1$ with geometric group
$\mathrm{PSL}_2(\F_l)$, then solvable covers of $X_0(l)$ and the
Atkin--Lehner quotient $X_0^*(l)$ of genus at most two.
Consequently the same prime-to-$p$ conclusion follows if additionally
$p\nmid l(l^2-1)$: the modular pullback, its Borel-subgroup steps,
and the order-two Atkin--Lehner step all have prime-to-$p$ Galois
closure. In the $l=5$ case, the new Galois map to $X(5)\simeq
\mathbf P^1$ has repeated branch indices and the third case applies.
Without the additional group-order check, tameness of the local
indices $2,3,l$ alone is insufficient to certify the required
Galois closure.\citep[Lemma~2.3 and \S3.4]{LA_Poonen2005}

\paragraph{The general reduction still required.}
B.LIT2 also applies whenever a prime-to-$p$ \'etale cover of the
given odd-characteristic curve maps to a hyperelliptic curve.
B.LIT4 establishes this under the listed low-genus Galois
hypotheses, but a general curve is not assumed to have such a map.

For a general curve, a tame map to $\mathbf P^1$ does not suffice.
Its ramification over a branch value may include unramified points,
where a branched target cover would remain ramified after base
change. Nor need the Galois closure of a tame map have
prime-to-$p$ order. Thus the tame Belyi theorem does not provide
the missing divisible fibres or the missing prime-to-$p$ covering
tower. The general universal-correspondence assertion is explicitly
Conjecture~2.6 in Bogomolov--Tschinkel; the theorem there is for
hyperelliptic sources and $p\ge5$. Neither that conjecture nor the
present special cases settle B5 for every nonhyperelliptic curve
or for all curves in characteristic two.
\citep[Theorem~1.7 and \S3]{LA_Poonen2005}
\citep[Theorem~1.7 and Conjecture~2.6]{LA_BogomolovTschinkel2002}
\citep[Theorem~1.5]{LA_KedlayaLittWitaszek2023}

The ordinarity literature does not supply the missing conclusion.
Raynaud's Theorem~2 constructs, for a proper smooth curve of genus
at least two in characteristic $p$, an \'etale cover with a solvable
prime-to-$p$ group which is nonordinary. Nonordinarity is weaker than
having a supersingular isogeny factor, and still weaker than the
required map to a supersingular curve. Thus this result cannot be
substituted for B5. The Prym results of Ozman--Pries likewise give
$p$-rank information, with no universal supersingular-target assertion.
\citep[Theorem~2]{LA_Raynaud2000}
\citep[Theorem~1.1 and \S8.2]{LA_OzmanPries2019}
